\appendix
\section{LLM Declaration}

\subsection{Tools used}

\begin{itemize}
    \item \textit{[Tool/version; access method; $22.10-9.11.2026$]}
    \item \textit{GPT 6-Luna (codex), $22.10-9.11.2026$}
    \item \textit{[GPT 6-Luna/version; GitHub Copilot (coding assistant session); $22.10-9.11.2026$]}
\end{itemize}
\subsection{Text contributions}

level 0 for no assistance.  
Level 1 for language-only edits
Level 2 for editorial changes and structural drafting
Level 3 for substantial drafting. 

\begin{description}
    \item[Abstract:] Level \textit{[3]}; \textit{[Copilot read through the rapport and the project, and wrote a summary which i verified and rewrote into the final abstract]}
    \item[Introduction:] Level \textit{[1]}; \textit{[Copilot created a basic motivation text which i used for inspiration, for the Motivation section.]}
    \item[Theory:] Level \textit{[2]}; \textit{[Copilot created the skeleton for equations, help with formatting and spacing of equations. Rewrote some equation descriptions]}
    \item[Backpropagation:] Level \textit{[2]}; \textit{[Copilot rewrote all equations describing backpropagation]}
    \item[Methods:] Level \textit{[2]}; \textit{[Copilot help with lists, figures and tables, formatting and structuring]}
    \item[Methods Subsections:] 
    \item[SGD:] Level \textit{[3]}; \textit{[ChatGPT-Sol 6.1 generated text describing the functions.]}
    \item[NNRegressor:] Level \textit{[3]}; \textit{[ChatGPT-Sol 6.1 generated text describing the functions.]} 
    \item[NNClassifier:] Level \textit{[3]}; \textit{[ChatGPT-Sol 6.1 generated text describing the functions.]}
    \item[Results:] Level \textit{[2]}; \textit{[Copilot help with lists, figures and tables, formatting and structuring]}
    \item[Discussion:] Level \textit{[0]}; \textit{[Describe assistance, or write ``none'']}
    \item[Conclusion:] Level \textit{[0]}; \textit{[Describe assistance, or write ``none'']}
\end{description}


\subsection{Code contributions}

List each code file or notebook included in the submission. Use level 0 for no
assistance, level 1 for debugging, level 2 for a snippet, level 3 for a
skeleton, and level 4 for substantial generation. Include what was tested,
adapted, or verified.
\begin{description}
    \item[\texttt{[Project2/Strucutre*]}:] Level \textit{[3]};
    \textit{[Copilot imported all neccesary classes and functions for the python files to function]}
\end{description}

\begin{description}
    \item[\texttt{[Project2/Code/Project2\_Classification\_Init.ipynb]}:] Level \textit{[2-3]};
    \textit{[Usage of classes objects. Translating functions into classes. Generated the one hot encoding of the target valeus, edited \texttt{predict\_proba} and predict to function properly. Editorial changes and shape errors in training. Generate Heatmaps and Display of data using GridSearchCV. Usage of and retrieving data from GridsearchCV]}
\end{description}

\begin{description}
    \item[\texttt{[Project2/Code/Project2\_Regression\_Init.ipynb]}:] Level \textit{[2-3]};
    \textit{[Usage pf classes and objects. Translating functions into classes. Debugging jax gradients and \texttt{compute\_gradients}. Creating activation functions inside NeuralNetRegressor. Editorial changes and shape errors in training. Generate Heatmaps and Display of data using GridSearchCV. Usage of and retrieving data from GridsearchCV. ]}
\end{description}

\subsubsection{LLM Declaration this section:}
The skeleton of this section was created using GitHub Copilot.  
(GPT Luna 6.0; coding-assistant session). 
